\begin{lemma}[A finite exterior-power test for monodromy purity]
\label{c:exterior-weight-criterion}
\status{Proved}
Let $V$ be a finite-dimensional mixed Frobenius module over
$\Q_\ell$, with integral weights and a nilpotent operator
$N:V\to V(-1)$.  Fix an integer $m$.  We use geometric Frobenius $\phi$ and
$\phi N\phi^{-1}=q^{-1}N$.  Suppose
\[
 \sum_w w\dim V_w=m\dim V,
\]
where $V=\bigoplus_wV_w$ is the decomposition by Frobenius
weights; Frobenius need not be semisimple.  Let $h$ be the
largest length of a Jordan string of $N$, with $h=0$ if $V=0$.
Then the following conditions are equivalent:
\begin{enumerate}
\item $V$ is pure of weight $m$ in the weight--monodromy sense;
\item for every $1\leq r\leq h$, all weights of
$\ker(N:\bigwedge^r V\to\bigwedge^r V(-1))$ are at most $rm$.
\end{enumerate}
Here the exterior power is the ordinary exterior power of vector
spaces, and $N$ acts on it as a derivation.  The determinant
hypothesis holds in particular when $V$ has a nondegenerate
Frobenius-equivariant pairing
$V\otimes V\to\Q_\ell(-m)$.
\end{lemma}

\begin{proof}
The generalized Frobenius eigenspaces of a given weight have a
sum defined over $\Q_\ell$: algebraic conjugation preserves the
weight of a Weil number.  They give the asserted direct
decomposition, and the Frobenius relation implies
$N(V_w)\subset V_{w-2}$.  None of these statements requires
Frobenius semisimplicity.

We first construct homogeneous Jordan strings.  Put
$K_L=\ker N\cap\operatorname{im}N^{L-1}$.  These are graded
subspaces with $K_{L+1}\subset K_L$.  Choose a homogeneous basis
$B_L$ of a complement of $K_{L+1}$ in $K_L$, and for every
$b\in B_L$ choose a homogeneous lift $v_b$ with
$N^{L-1}v_b=b$.  The strings
\[
 v_b,Nv_b,\ldots,N^{L-1}v_b
\]
are linearly independent together.  Indeed, in a relation,
apply the largest power of $N$ equal to the distance of a
remaining term from its bottom.  The surviving bottom vectors
belong to the independent union of the $B_L$, so their
coefficients vanish; descending induction finishes the argument.
They have the correct total dimension, since
\[
 \dim K_L=\operatorname{rank}N^{L-1}
               -\operatorname{rank}N^L,\qquad
 \sum_L L(\dim K_L-\dim K_{L+1})=\dim V.
\]
Thus they form a homogeneous Jordan basis.  The individual
strings need not be Frobenius-stable.

A string of length $L$, whose top vector has weight $b$, has
weights $b,b-2,\ldots,b-2(L-1)$ and centre $c=b-L+1$.
The monodromy filtration assigns these positions the indices
$L-1,L-3,\ldots,-L+1$.  Hence weight--monodromy purity of
weight $m$ is equivalent to every string centre being $m$.

The full wedge of the string,
\[
 v_b\wedge Nv_b\wedge\cdots\wedge N^{L-1}v_b,
\]
is nonzero, is killed by $N$, and has weight
\[
 Lb-2(0+\cdots+L-1)=L(b-L+1)=Lc.
\]
Each term in its derivative has a repeated vector or is zero.
Although the wedge line need not be Frobenius-stable, the
nonzero intersection of $\ker N$ with the weight-$Lc$ subspace
is Frobenius-stable.  Condition (2), applied with $r=L$, therefore
gives $c\leq m$ for every string.  The determinant hypothesis is
$\sum_{\text{strings}}L(c-m)=0$, so every centre is $m$.

Conversely, if every centre is $m$, the homogeneous strings
carry the usual $\mathfrak{sl}_2$ structure whose lowering
operator is $N$ and whose Cartan weights are the Frobenius
weights minus $m$.  On $\bigwedge^rV$ the Cartan weights are
the Frobenius weights minus $rm$.  A vector $z$ of Cartan
weight $a$ killed by the lowering operator has $a\leq0$:
if $t$ is maximal with $e^tz\ne0$, the relation
$f e^kz=-k(a+k-1)e^{k-1}z$ at $k=t+1$ gives
$a=-t$.  Here $e$ is nilpotent because it raises the finite
weight grading.  This proves (2), in fact for every exterior power.
Finally, a perfect pairing of the indicated weight pairs
$V_w$ perfectly with $V_{2m-w}$, and gives the determinant
hypothesis by summing dimensions.

This is a direct graded-string argument.  The related
determinant criterion and invariance of purity under finite
extension are recorded, with their proofs, in
\citet[Lemma~1.4(1), (2), (5), author pp.~6--7]{c_tayloryoshida2007}.
\end{proof}

\begin{theorem}[The universal trait endpoint and weight--monodromy]
\label{c:universal-trait-weight-equivalence}
\status{Proved equivalence; the universal assertions remain unresolved}
Fix a residue characteristic $p$ and a prime $\ell\ne p$.
We consider finite extensions $K/\Q_p$, with finite residue
field, and geometric $\ell$-adic cohomology.  The following
universal assertions are equivalent:
\begin{enumerate}
\item Weight--monodromy holds for $H^j(Z_{\bar K},\Q_\ell)$
for every smooth projective $Z/K$ and every $j\geq0$.
\item For every smooth projective $j$-fold $Z/K$,
$H^j(Z_{\bar K},\Q_\ell)^{I_K}$ has weights at most $j$.
\item The assertion in (2) holds for the generic fibres of
projective strictly semistable models over $\mathcal O_K$.
\end{enumerate}
In (2) the stated space is the closed geometric stalk of
$u_*H^j(Z_{\bar K},\Q_\ell)$ for the open immersion
$u:\Spec K\hookrightarrow\Spec\mathcal O_K$.  Thus it is exactly
the dimension-normalized, degree-zero trait endpoint.

For an individual $j$-fold, after a finite extension making its
inertia action unipotent, it suffices to test the bound
$rj$ on $\ker N$ in
$\bigwedge^r H^j(Z_{\bar K},\Q_\ell)$ for $1\leq r\leq j+1$.
Each such test is a direct-summand test in the middle cohomology
$H^{rj}(Z^r_{\bar K},\Q_\ell)$.  A failure of weight--monodromy
therefore yields a trait-endpoint failure on one of these powers,
or on a projective semistable alteration of that power.
\end{theorem}

\begin{proof}
We give the transfers explicitly.  Finite extension multiplies
$N$ by a nonzero scalar if the tame character is renormalized
and replaces Frobenius weights by the same weights relative to
the larger residue cardinality.  For the latter assertion one
may pass to a common normal extension killing the finite
Weil--Deligne inertia, and choose a positive power of a
Frobenius lift lying in its Weil group.  A further power of the
new Frobenius lift differs from it by inertia acting trivially
in the Weil--Deligne representation.  Taking roots of the
corresponding absolute values proves the assertion.
Consequently weight--monodromy is invariant under finite
extension; this also verifies directly
\citet[Lemma~1.4(2)]{c_tayloryoshida2007}.

For a projective semistable model, the spectral sequence used
in Proposition~\ref{c:improved-zero-weight} has
\[
 E_1^{a,b}=
 \bigoplus_{s\geq\max(0,-a)}
 H^{b-2s}(Y^{(a+2s)}_{\bar k},\Q_\ell)(-s).
\]
Here $Y^{(t)}$ denotes the disjoint union of the $(t+1)$-fold
intersections.  On total degree $j=a+b$, a nonzero summand
requires
$b\geq2s\geq\max(0,2b-2j)$, hence $0\leq b\leq2j$.
The terms have weight $b$, so all abutment weights lie in
$[0,2j]$.  Since $N$ lowers weights by two,
$N^{j+1}=0$.  The spectral-sequence and monodromy statements
are precisely \citet[Satz~2.10 and (2.11), pp.~40--41]{c_rapoportzink1982}.
For arbitrary smooth projective $Z$, the same conclusions
follow after a projective semistable alteration and finite
extension: pullback is split by pushforward divided by the
generic degree, so its cohomology is a Frobenius- and
$N$-stable direct summand.  This is the alteration and trace
reduction already checked in Proposition~\ref{c:improved-zero-weight},
using \citet[Theorem~6.5, pp.~83--86]{c_dejong1996}.

For a $j$-fold, Poincare duality pairs $H^j$ perfectly into
$\Q_\ell(-j)$.  Lemma~\ref{c:exterior-weight-criterion}
therefore applies with largest string length at most $j+1$.
The required ordinary exterior power is a Galois direct
summand of the Kunneth factor
$H^j(Z)^{\otimes r}\subset H^{rj}(Z^r)$.  Explicitly its
projector on that factor is
\[
 \frac1{r!}\sum_{\sigma\in\mathfrak S_r}
       (\operatorname{sgn}\sigma)^{j+1}\sigma^*.
\]
The geometric permutation $\sigma^*$ already contributes
the Koszul factor $(\operatorname{sgn}\sigma)^j$.
Thus we use the geometric antisymmetrizer for even $j$ and
the geometric symmetrizer for odd $j$.  These are direct
summands of Galois representations; no algebraic
Chow--Kunneth projector is assumed.  The variety $Z^r$ has
dimension $rj$, exactly as required by the normalization.
The rational denominator $r!$ causes no exception at $\ell=2$.

Assume (3).  Given a $j$-fold $Z$, pass first to a finite
extension on which its inertia action is unipotent.  For each
$1\leq r\leq j+1$, choose, after a further finite extension,
a projective strictly semistable alteration
$\pi_r:W_r\to Z^r$ on generic fibres over some $L_r$.
We allow the fields $L_r$ to depend on $r$; no ramified
base change of a semistable integral model is presumed to
remain regular.
The maps
\[
 H^{rj}(Z^r)\xrightarrow{\pi_r^*}H^{rj}(W_r)
 \xrightarrow{(\deg\pi_r)^{-1}\pi_{r*}}H^{rj}(Z^r)
\]
have composite the identity and commute with inertia and
monodromy.  On a semistable representation, inertia invariants
are exactly $\ker N$, because its inertia action is
$\exp(t_\ell(\tau)N)$ and $t_\ell$ has infinite image.
Assumption (3) consequently bounds the weights of
$\ker N$ on each indicated exterior summand over $L_r$ by $rj$.
The original inertia action was already unipotent, so passage
to $L_r$ leaves this kernel unchanged.  Frobenius-power
comparison descends the same bound separately for every $r$.
The lemma proves weight--monodromy for $H^j(Z)$.
Finite extension does not change that conclusion.

We next reduce arbitrary cohomological degrees to this middle
case, retaining actual split summands.  Let $\dim X=n$.
If $n\geq j$, choose a smooth complete intersection
$i:Z\hookrightarrow X$ of dimension $j$, for a very ample
class $L$.  Weak Lefschetz and hard Lefschetz give the
Galois-equivariant retraction
\[
 H^j(Z)\xrightarrow{i_*}
 H^{2n-j}(X)(n-j)
 \xrightarrow{(L^{n-j})^{-1}}H^j(X),
\]
whose composite with $i^*$ is the identity.  The inverse here
is an inverse of cohomological Galois modules; its algebraicity
is unnecessary.  If $n<j$, replace $X$ by
$X\times\mathbf P^{j-n}$ and take the $H^0$ Kunneth factor
of projective space, again an actual split summand.
Pure weight--monodromy modules are closed under split
submodules: the canonical monodromy filtration is additive
on a direct sum, so each summand inherits the pure graded
pieces.  Thus (3) implies (1).

If (1) holds, a monodromy string of centre $j$ has bottom
weight $j-L+1\leq j$.  After a semistable extension this proves
the bound on $\ker N$.  The original $I_K$-invariants form
a subspace of these invariants; comparison by Frobenius
powers, on which the original inertia acts trivially, gives
the same upper bound over $K$.  Hence (1) implies (2).
Finally (2) implies (3) by restriction of its range.
\end{proof}

\begin{example}[One exterior test can miss a failure]
\label{c:single-endpoint-insufficient}
\status{Proved formal counterexample; no geometric realization asserted}
Take $K=\Q_2$, $\ell=3$, $q=2$, and an unramified Weil
representation over $\Q_3$ with basis
$u_6,v_4,x_4,y_2$, subscripts denoting weights.  Put
\[
 \phi=\operatorname{diag}(8,4,4,2),\qquad
 N(u_6)=v_4,\quad N(x_4)=y_2,\quad N(v_4)=N(y_2)=0.
\]
Its two string centres are $5$ and $3$, so it is not
weight--monodromy pure of weight $4$.  Nevertheless
its determinant has weight $16=4\dim V$ and
$\ker N=\langle v_4,y_2\rangle$ has weights at most $4$.
It even has the nondegenerate symmetric pairing into
$\Q_\ell(-4)$ specified by
\[
 \langle u_6,y_2\rangle=1,\qquad
 \langle v_4,x_4\rangle=-1
\]
and symmetric counterparts, all other entries being zero.
The identities $\phi N\phi^{-1}=2^{-1}N$,
$N^{\mathsf T}Q+QN=0$, and
$\phi^{\mathsf T}Q\phi=2^4Q$ are immediate matrix identities.
In $\bigwedge^2V$, the nonzero vector $u_6\wedge v_4$
is killed by $N$ and has weight $10>8$.
The second exterior test detects the failure.
The exact rational-matrix certificate
\path{code/part_c/check_c4_exterior_witness.py}
and its companion
\path{code/part_c/check_c4_exterior_witness_output.json}
verify these identities and the exterior-power calculation.
This is a counterexample only to the proposed linear-algebra
implication from the single invariant bound and self-duality.
It is not a smooth projective fourfold or a counterexample
to geometric weight--monodromy.
\end{example}

\begin{remark}[Scope of the trait criterion]
\label{c:trait-exterior-scope}
\status{Proved scope clarification; universal C.4 remains incomplete}
The theorem identifies the universal middle-degree trait
endpoint with weight--monodromy and gives a finite power
witness to a failure.  It does not prove either universal
assertion.  Its bound $r\leq j+1$ uses the actual nonnegative
cohomological degree $j$; it is not a bound in terms of the
weight of an arbitrary Tate-twisted motive.

The finite inertia quotient can remove a vector from the
invariants at the original place, which is why the witness
is taken after a finite semistable extension.  All local fields
are quantified in the universal assertion, so this extension
is permitted.  The argument concerns the stated traits.
For a general morphism in C.4, the component and curve-selection
argument in Proposition~\ref{c:improved-zero-weight} is an
additional required transfer; no general sheaf bound is
inferred merely by naming the trait equivalence.
\end{remark}
